% !TeX spellcheck = en_US
% !TeX encoding = UTF-8
\chapter*{Abstract}

In the second lab of the Security Insider Lab, we detect network attacks from flow features with machine learning. We first recorded our own traffic and converted it into flows with CICFlowMeter. Our first captures were mostly IPv6 and produced bogus flows, because CICFlowMeter force-parses non-IPv4 frames as IPv4. With IPv6 disabled, we obtained 569 flows, 547 of them valid TCP/UDP flows. As training data we chose the Thursday web-attack file of CIC-IDS2017 (Brute Force, XSS, SQL injection) instead of the suggested DDoS file, because web attacks look like normal HTTP and are harder to detect. Pre-processing removed infinite values, duplicate rows and constant and identical columns, which left 164,179 flows with 60 features. We then made a stratified split, min-max scaled the features and selected the 20 best with $\chi^2$. We trained Logistic Regression and Decision Tree classifiers for a binary and a four-class target. The default Decision Tree is best: it detects all 643 attacks of the test split with 18 false positives (attack F1 0.986) and reaches a multiclass macro-F1 of 0.729, whereas default multiclass Logistic Regression reaches 98.5\,\% accuracy while detecting no attack. Applied flow by flow to our own benign capture, the Decision Tree raised 3 false alarms in 547 flows (0.55\,\%, about 15 times the test-split rate), caused by shortcut features such as the destination port and the initial TCP window size.
